\section*{Chapter \ref{ch:nw:programmable-hardware-ii-trading-designs} --- Programmable Hardware II: Trading Designs}

\begin{solution}{exo:nw:programmable-hardware-ii-trading-designs:1}
The price ends at byte $20 + 2 + 35 = 57$, in beat 7 (bytes 56 to 63). The order is registered two cycles later, at the end of cycle 9:
ten cycles from the start of the packet's first beat, $10 \times 6.4 = \qty{64}{\nano\second}$.
\end{solution}

\begin{solution}{exo:nw:programmable-hardware-ii-trading-designs:2}
One token every 64 cycles: $156.25 \times 10^6 / 64 \approx 2.44$ million orders a second sustained, after an initial burst of 4.
\end{solution}

\begin{solution}{exo:nw:programmable-hardware-ii-trading-designs:3}
$50 \times 156.25 = 7\,812.5$ cycles; $4 + \lfloor 7\,812.5 / 64\rfloor = 4 + 122 = 126$ orders.
\end{solution}

\begin{solution}{exo:nw:programmable-hardware-ii-trading-designs:4}
The 47 extra qualifying messages (160 against 113) arrive in bursts, when the offers near the threshold are being placed: the bucket, already
drained by the first orders of each burst, rejects 22 of them and lets 25 through.
\end{solution}

\begin{solution}{exo:nw:programmable-hardware-ii-trading-designs:5}
An order at the limit is not the order the strategy computed; clipping silently changes risk and hides a configuration or model error.
Rejecting it keeps the hardware's behaviour simple and auditable and makes the error visible in the reject counter.
\end{solution}

\begin{solution}{exo:nw:programmable-hardware-ii-trading-designs:6}
The threshold itself (a model over many inputs and history), the position and exposure across venues (state shared with other systems),
and whether to trade at all (news, halts, the desk's limits). Each needs state, arithmetic or judgement that would cost logic and weeks of
verification for every change.
\end{solution}

\begin{solution}{exo:nw:programmable-hardware-ii-trading-designs:7}
One order: the fixture lasts fewer than 10\,000 cycles, so the bucket never refills after its single token. With the kill bit set one cycle
before the tenth order of the unrestricted run, exactly nine orders leave.
\end{solution}

\begin{solution}{exo:nw:programmable-hardware-ii-trading-designs:8}
The 13 nanoseconds (a benchmark's figure) runs from the last bit needed for the decision to the first bit of the order. Wire to wire adds the
time for the triggering bits to arrive, the card's transceivers, and the cage's fibre and devices both ways: hundreds of nanoseconds in
the chapter's cage.
\end{solution}

\begin{solution}{pb:nw:programmable-hardware-ii-trading-designs:1}
\begin{enumerate}
\item $263 + 120 + 19.2 + 120 + 295 \approx \qty{817}{\nano\second}$: \qty{0.82}{\micro\second}.
\item Three cycles, \qty{19.2}{\nano\second}: about 2\%.
\item \qty{0.82}{\micro\second}: every stage of the model is a constant (fibre, fixed transceiver latency, fixed cycles).
\item By \qty{3.57}{\micro\second} at the median and \qty{34.5}{\micro\second} at the 99th percentile.
\item 75\% for the hardware path, 0.01\% for the software path.
\item 100\% and 57\%.
\item The cage's fibre and the transceivers: what is left is physics (distance) and serialisation, chapters~9 to~14.
\item The competitor's correlation with the firm (both see the same bursts), queueing at the venue's gateway, and anything but a single race.
\item $4 + \lfloor 3\,125/64\rfloor = 52$ orders.
\item $4 + \lfloor 16/64\rfloor = 4$ orders: the bucket's burst.
\item So that an order already past the risk stage is stopped in its last cycle too: the kill bit then blocks everything from the cycle it
  is set.
\item On the card: size, price collar against a reference, order and notional rates, the kill bit. In the broker's layer, under its control:
  credit and capital thresholds, and its own kill switch.
\item \textbf{Named result.} \qty{0.82}{\micro\second} wire to wire in hardware against a median of \qty{4.39}{\micro\second} in software; against a
  competitor at \qty{1}{\micro\second} they win 75\% and 0.01\% of races; a \qty{20}{\micro\second} kill lets at most 52 orders out.
\item Hardly: at \qty{5}{\micro\second} the software path already wins 57\% and the rest is spread across its tail; better to fix the tail.
\item Flexibility, speed of change and much of the strategy's expressiveness; it gains a verification burden.
\item As often as software can compute and write it, within microseconds; the strategy's software owns it, under the risk layer's limits.
\item To know its position and to reconcile with the venue's acknowledgements: the card acts alone, and the firm must not lose track of
  what it did.
\item Replay the recorded feed through the HDL in simulation and through the card in a test harness, compare both with the models, and run
  the venue's certification.
\item Stale thresholds or limits: the design should hold a version number and a heartbeat from software, and stop firing when they are
  stale.
\item Hardware decides when; software decides what, and whether at all.
\end{enumerate}
\end{solution}

\begin{solution}{iq:nw:programmable-hardware-ii-trading-designs:1}
It processes a fixed number of bytes per cycle with a state machine that tracks the position in the header and in each length-prefixed
message, capturing fields as their bytes pass; a message's fields are known the cycle its last byte arrives, and decisions can start
before the packet ends or is checked.
\iqlookfor{streaming state, not buffering, and deciding before the frame ends.}
\end{solution}

\begin{solution}{iq:nw:programmable-hardware-ii-trading-designs:2}
Size, price collar, rate limits (a token bucket), duplicate detection and a kill bit, each a comparison or counter in a pipeline stage that
every order must cross. Prove it by construction (the only path to the output goes through the stage), by simulation of every boundary, and
by keeping the broker's controls independent of the firm's configuration.
\iqlookfor{a single path to the wire and independent controls.}
\end{solution}

\begin{solution}{iq:nw:programmable-hardware-ii-trading-designs:3}
Models, parameters, positions, cross-venue risk and anything needing history or judgement: software computes what would be worth doing,
hardware applies it at the moment the market offers it. Hardware changes are slow and costly to verify; software changes are cheap.
\iqlookfor{what against when, and the cost of change.}
\end{solution}

\begin{solution}{iq:nw:programmable-hardware-ii-trading-designs:4}
With a \omterm{def:nw:programmable-hardware-ii-trading-designs:toe}{TCP offload engine} on the card that owns the connection's sequence numbers, acknowledgements and retransmissions, and an \omterm{def:nw:programmable-hardware-ii-trading-designs:template}{order
template} whose TCP and IP headers and checksums are patched at send time; the session layer (logon, heartbeats) is shared with software.
\iqlookfor{connection state in hardware and a template with patched fields.}
\end{solution}

\begin{solution}{iq:nw:programmable-hardware-ii-trading-designs:5}
Run both on the same recorded and synthetic inputs, in two simulators, and compare every output with its cycle; add message-level reference
checks (each qualifying message decided once). A disagreement is a bug in one of them until shown otherwise: reduce the input to the
smallest case and read the waveform.
\iqlookfor{cycle-level comparison and disciplined reduction.}
\end{solution}

\begin{solution}{iq:nw:programmable-hardware-ii-trading-designs:6}
From which bit to which bit it is measured, on which frame sizes and message types, whether it includes the transceivers and the order's
TCP, under what load and with what risk checks enabled, and how the figure maps to my wire-to-wire path.
\iqlookfor{the definition of the measurement.}
\end{solution}
